% =====================================================================
%  5.22 ANEXOS (OPCIONALES)
%  No cuentan para el limite de 50 paginas: es el sitio para todo el
%  material voluminoso que en el cuerpo estorbaria.
%  main.tex ya ejecuta \appendix antes de este input.
% =====================================================================

\chapter{Comandos de operación}
\label{anexo:comandos}

Este anexo resume los comandos necesarios para instalar, ejecutar,
comprobar y auditar el prototipo desde la raíz del repositorio. El
recorrido de uso se describe en el~\cref{cap:manual}; aquí solo se
documenta la operación técnica. El servidor de desarrollo debe lanzarse
con la configuración \texttt{development}; no se emplea la
compilación de producción como entorno de trabajo.

\begin{table}[H]
  \centering
  \small
  \begin{tabularx}{\textwidth}{@{}>{\raggedright\ttfamily}p{7.4cm}X@{}}
    \toprule
    \textbf{\normalfont Comando} & \textbf{Propósito} \\
    \midrule
    npm install &
      Instala las dependencias declaradas en \texttt{package.json}. \\
    npx ng serve
    --host 127.0.0.1 --port 4200 &
      Arranca el servidor de desarrollo (configuración
      \texttt{development}) en el puerto 4200. \\
    npx ng test --watch=false
    --browsers=ChromeHeadless &
      Ejecuta la batería Karma + Jasmine en Chrome sin interfaz, una
      sola pasada. \\
    npx ng build
    --configuration development &
      Genera el artefacto de desarrollo (comprobación de tipos incluida). \\
    bash scripts/check-links.sh &
      Verifica el patrón de documentación enlazada (\texttt{@linked}):
      enlaces rotos, documentos huérfanos y ficheros fuera del mapa. \\
    node scripts/audit-criteria.cjs &
      Audita la consistencia entre \texttt{criteria.json}, el catálogo
      de fármacos y el de diagnósticos (clases y códigos huérfanos). \\
    \bottomrule
  \end{tabularx}
  \caption{Comandos de instalación, ejecución, prueba y auditoría.}
  \label{tab:comandos-operacion}
\end{table}

El script \texttt{scripts/check-links.sh} existe en el repositorio y
debe terminar con código de salida 0. Las líneas de consola del tipo
«Error evaluando criterio» que puedan aparecer durante
\texttt{ng test} corresponden a una especificación deliberada de manejo
de errores, no a un fallo de la suite: el veredicto es la línea final
de totales.

\chapter{Sistemas clínicos y recuento de criterios}
\label{anexo:sistemas}

La guía STOPP/START v3 formula 190 enunciados (133 STOPP y 57
START)~\cite{omahony2023stopp}. El catálogo de la aplicación no los
reproduce 1:1: algunos enunciados se desglosan en varias reglas
(subcriterios) y otros se agrupan. La~\cref{tab:sistemas-criterios}
resume, por los 13 sistemas del fichero \texttt{criteria.json}, cuántas
reglas STOPP y START hay implementadas. No se listan los textos
completos: el volumen duplicaría la memoria sin añadir información de
diseño.

Los recuentos corresponden a la versión del catálogo presente en el corte
documental del 16 de agosto de 2026 (218 reglas: 166 STOPP y 52 START). Son cifras de
implementación, no un censo de la guía impresa~\cite{sacyl2023stopp}.

\begin{table}[H]
  \centering
  \small
  \begin{tabular}{@{}lrrr@{}}
    \toprule
    \textbf{Sistema clínico} & \textbf{STOPP} & \textbf{START} & \textbf{Total} \\
    \midrule
    Indicación de la medicación                    &   8 &  0 &   8 \\
    Sistema cardiovascular                         &  35 & 11 &  46 \\
    Anticoagulantes/Antiagregantes                 &  19 &  2 &  21 \\
    Sistema nervioso central                       &  31 &  7 &  38 \\
    Sistema renal                                  &  10 &  4 &  14 \\
    Sistema gastrointestinal                       &   9 &  7 &  16 \\
    Sistema respiratorio                           &   4 &  3 &   7 \\
    Sistema musculoesquelético                     &  11 &  9 &  20 \\
    Sistema urogenital                             &   8 &  5 &  13 \\
    Sistema endocrino                              &  11 &  1 &  12 \\
    Riesgo de caídas                               &  14 &  0 &  14 \\
    Analgésicos                                    &   5 &  3 &   8 \\
    Carga antimuscarínica/anticolinérgica          &   1 &  0 &   1 \\
    \midrule
    \textbf{Total}                                 & \textbf{166} & \textbf{52} & \textbf{218} \\
    \bottomrule
  \end{tabular}
  \caption{Reglas implementadas por sistema clínico (catálogo
    \texttt{criteria.json}).}
  \label{tab:sistemas-criterios}
\end{table}

\chapter{Trazabilidad entre requisitos y evidencia}
\label{anexo:trazabilidad}

La~\cref{tab:trazabilidad-requisitos} enlaza cada requisito con evidencia
localizada en el repositorio. «Localizada» significa que existe una
especificación automática que expresa ese comportamiento; no significa que
la suite se haya ejecutado con éxito para esta edición de la memoria.
«Parcial» identifica requisitos cuya dimensión completa exige medición,
auditoría o prueba de uso adicional. Salvo las rutas que comienzan por
\texttt{src/}, los nombres de fichero se refieren a \texttt{src/app/}.

{\scriptsize
\begin{longtable}{@{}p{0.7cm}p{3.4cm}p{6.0cm}p{2.7cm}@{}}
  \caption{Matriz requisitos--evidencia del repositorio.}
  \label{tab:trazabilidad-requisitos} \\
  \toprule
  \textbf{ID} & \textbf{Comportamiento trazado} &
  \textbf{Evidencia localizada} & \textbf{Estado documental} \\
  \midrule
  \endfirsthead
  \caption[]{Matriz requisitos--evidencia del repositorio (cont.).} \\
  \toprule
  \textbf{ID} & \textbf{Comportamiento trazado} &
  \textbf{Evidencia localizada} & \textbf{Estado documental} \\
  \midrule
  \endhead
  \bottomrule
  \endfoot
  RF01 & Seleccionar y deseleccionar medicamentos &
    \texttt{steps/meds-step/meds-step.component.spec.ts};
    \texttt{core/group-checked.spec.ts} &
    Localizada; ejecución final no registrada \\
  RF02 & Seleccionar diagnósticos y variantes excluyentes &
    \texttt{steps/diagnosis-step/diagnosis-step.component.spec.ts};
    \texttt{core/data/diagnosis-variants.spec.ts} &
    Localizada; ejecución final no registrada \\
  RF03 & Capturar analítica y constantes &
    \texttt{core/lab-capture.spec.ts};
    \texttt{steps/diagnosis-step/diagnosis-step.component.spec.ts} &
    Localizada; ejecución final no registrada \\
  RF04 & Reevaluar tras cambios del caso &
    \texttt{core/services/criteria-engine.service.spec.ts};
    especificaciones de ambos pasos &
    Localizada; ejecución final no registrada \\
  RF05 & Calcular advertencias preventivas de medicación &
    \texttt{core/services/criteria-engine.service.spec.ts} &
    Solo motor; no integrado en la interfaz; ejecución final no registrada \\
  RF06 & Mostrar elementos relevantes de otros sistemas &
    \texttt{core/data/system-relevance.spec.ts};
    \texttt{core/group-visibility.spec.ts} &
    Localizada; ejecución final no registrada \\
  RF07 & Marcar pestañas revisadas y persistir el estado &
    \texttt{core/case-store.service.spec.ts} &
    Localizada; ejecución final no registrada \\
  RF08 & Agrupar resultados por tipo y sistema &
    \texttt{core/criteria-groups.spec.ts};
    especificaciones de ambos pasos &
    Localizada; ejecución final no registrada \\
  RF09 & Generar y descargar el informe PDF &
    \texttt{core/report.service.spec.ts} &
    Localizada; no se adjunta PDF ejecutado \\
  RF10 & Copiar criterios como texto plano &
    \texttt{core/clipboard-text.spec.ts};
    especificaciones de ambos pasos &
    Localizada; ejecución final no registrada \\
  RF11 & Exportar e importar JSON validado &
    \texttt{core/case-io.service.spec.ts} &
    Localizada; no se adjunta JSON ejecutado \\
  RF12 & Reiniciar solo tras confirmación &
    \texttt{confirm-reset-dialog.component.spec.ts};
    \texttt{core/case-store.service.spec.ts} &
    Localizada; ejecución final no registrada \\
  RF13 & Escala de fuente y orientación persistentes &
    \texttt{core/display-settings.service.spec.ts};
    especificaciones de ambos pasos &
    Localizada; ejecución final no registrada \\
  RF14 & Guía rápida y enlace a la guía oficial &
    \texttt{quick-guide-dialog.component.spec.ts} &
    Localizada; ejecución final no registrada \\
  RF15 & Opción «Otro» por grupo &
    \texttt{core/group-checked.spec.ts};
    especificaciones de ambos pasos &
    Localizada; ejecución final no registrada \\
  RNF01 & Flujo en dos pasos con resultados junto a la captura &
    \texttt{app.routes.spec.ts}; especificaciones de ambos pasos &
    Parcial: estructura localizada; usabilidad no evaluada \\
  RNF02 & Evaluación en cliente en la misma interacción &
    \texttt{core/services/criteria-engine.service.spec.ts};
    especificaciones de ambos pasos &
    Parcial: sin medición de latencia \\
  RNF03 & Tres escalas de fuente &
    \texttt{core/display-settings.service.spec.ts} &
    Parcial: sin auditoría de accesibilidad \\
  RNF04 & Catálogo JSON separado del motor &
    \texttt{src/assets/data/criteria.json};
    \texttt{core/data/criteria-data-integrity.spec.ts} &
    Estructura localizada; ejecución final no registrada \\
  RNF05 & Persistencia local del caso &
    \texttt{core/case-store.service.spec.ts};
    \texttt{core/case-io.service.spec.ts} &
    Parcial: persistencia localizada; privacidad no auditada \\
\end{longtable}
}

\chapter{Cobertura demostrable del catálogo de criterios}
\label{anexo:cobertura-criterios}

El catálogo contiene 218 reglas JSON, distribuidas como muestra la
\cref{tab:sistemas-criterios}. Ese recuento se obtiene directamente de
\texttt{src/assets/data/criteria.json}. Las especificaciones
\texttt{criteria-a.spec.ts} a \texttt{criteria-m.spec.ts} contienen casos
por secciones del catálogo, y
\texttt{criteria-data-integrity.spec.ts} expresa guardas entre reglas y
catálogos.

La~\cref{tab:cobertura-catalogo} es deliberadamente una matriz de
\emph{evidencia del repositorio}, no una matriz de equivalencia clínica.
La presencia de un fichero por sección tampoco demuestra por sí sola que
cada una de sus reglas tenga casos positivo y negativo ejecutados.
No se ha validado una correspondencia criterio por criterio entre los 190
enunciados de la guía (133 STOPP y 57 START) y las 218 reglas. Una regla
puede desglosar parte de un enunciado, varias reglas pueden agruparse o un
criterio puede necesitar datos no capturados. Por tanto, ningún recuento de
esta tabla demuestra cobertura completa de la guía.

{\scriptsize
\begin{longtable}{@{}p{4.5cm}rrp{3.0cm}p{3.0cm}@{}}
  \caption{Matriz de cobertura demostrable del catálogo, sin equivalencia
  190$\rightarrow$218.}
  \label{tab:cobertura-catalogo} \\
  \toprule
  \textbf{Sistema del catálogo} & \textbf{STOPP} & \textbf{START} &
  \textbf{Especificación localizada} & \textbf{Correspondencia con guía} \\
  \midrule
  \endfirsthead
  \caption[]{Matriz de cobertura demostrable del catálogo (cont.).} \\
  \toprule
  \textbf{Sistema del catálogo} & \textbf{STOPP} & \textbf{START} &
  \textbf{Especificación localizada} & \textbf{Correspondencia con guía} \\
  \midrule
  \endhead
  \bottomrule
  \endfoot
  Indicación de la medicación & 8 & 0 & \texttt{criteria-a.spec.ts} & No validada \\
  Sistema cardiovascular & 35 & 11 & \texttt{criteria-b.spec.ts} & No validada \\
  Anticoagulantes/Antiagregantes & 19 & 2 & \texttt{criteria-c.spec.ts} & No validada \\
  Sistema nervioso central & 31 & 7 & \texttt{criteria-d.spec.ts} & No validada \\
  Sistema renal & 10 & 4 & \texttt{criteria-e.spec.ts} & No validada \\
  Sistema gastrointestinal & 9 & 7 & \texttt{criteria-f.spec.ts} & No validada \\
  Sistema respiratorio & 4 & 3 & \texttt{criteria-g.spec.ts} & No validada \\
  Sistema musculoesquelético & 11 & 9 & \texttt{criteria-h.spec.ts} & No validada \\
  Sistema urogenital & 8 & 5 & \texttt{criteria-i.spec.ts} & No validada \\
  Sistema endocrino & 11 & 1 & \texttt{criteria-j.spec.ts} & No validada \\
  Riesgo de caídas & 14 & 0 & \texttt{criteria-k.spec.ts} & No validada \\
  Analgésicos & 5 & 3 & \texttt{criteria-l.spec.ts} & No validada \\
  Carga antimuscarínica/anticolinérgica & 1 & 0 &
    \texttt{criteria-m.spec.ts} & No validada \\
  \midrule
  \textbf{Total del catálogo} & \textbf{166} & \textbf{52} &
    \textbf{Especificaciones localizadas; ejecución final no registrada} &
    \textbf{No validada} \\
\end{longtable}
}

Para convertir esta matriz en una acreditación de cobertura clínica falta
una revisión independiente que enumere los 190 criterios oficiales y, para
cada uno, identifique reglas, datos necesarios, caso positivo, caso
negativo y estado (completo, parcial, manual, no evaluable o no
implementado). Esa correspondencia no puede inferirse de los totales y no
se inventa en esta memoria.
